\subsection{Regular and Singular Cardinals}

\newcommand\cofin[1]{\mathrm{cf}(#1)}

\exercise{1}{
    $\cofin{\al_\w} = \cofin{\al_{\w + \w}} = \w$.
}
\sol{
    First we show that $\cofin{\al_\w} = \w$.
    \qproof{
        By definition,
        \gath{
            \al_\w = \sup\braces{\al_n \where n < \w} = \lim_{n \to \w} \al_n
        }
        so that $\cofin{\al_\w} \leq \w$.
        Moreover, we know that $\cofin{\a}$ is always a limit ordinal, and that $\w$ is the smallest nonzero limit ordinal.
        As there are no ordinals less than zero, we clearly have
        \gath{
            \lim_{n \to 0} \a_n = \sup \braces{\a_n \where n < 0} = \sup \es = \bigcup \es = \es \neq \al_w
        }
        so that it must in fact be that $\cofin{\al_\w} = \w$ as desired.
    }

    Now we show that $\cofin{\al_{\w + \w}} = \w$ as well.
    \qproof{
        We we clearly have
        \gath{
            \al_{\w + \w} = \lim_{n \to \w} \al_{\w + n},
        }
        noting that this is remarked upon following Theorem~9.2.2.
        Hence, $\cofin{\al_{\w + \w}} \leq \w$ and, by the same argument as above it has to be that $\cofin{\al_{\w + \w}} = \w$ since $\w$ is the smallest nonzero limit ordinal.
    }
}

\exercise{2}{
    $\cofin{\al_{\w_1}} = \w_1$, $\cofin{\al_{\w_2}} = \w_2$.
}
\sol{
    TODO
}

\exercise{3}{
    Let $\a$ be the cardinal number defined in the proof of Lemma~9.2.6.
    Show that $\cofin{\a} = \w$.
}
\sol{
    TODO
}

\exercise{4}{
    Show that $\cofin{\a}$ is the least $\g$ such that $\a$ is the union of $\g$ sets of cardinality less than $\abs{\a}$.
}
\sol{
    TODO
}

\exercise{5}{
    Let $\al_\a$ be a limit cardinal, $\a > 0$.
    Show that there is an increasing sequence of \emph{alephs} of length $\cofin{\al_\a}$ with limit $\al_\a$.
}
\sol{
    TODO
}

\exercise{6}{
    Let $\k$ be a limit cardinal, and let $\l < \k$ be a regular infinite cardinal.
    Show that there in an increasing sequence $\angles{\a_\n \where \n < \cofin{\k}}$ of cardinals such that $\lim_{\n \to \cofin{\k}} \a_\n = \k$ and $\cofin{\a_\n} = \l$ for all $\n$.
}
\sol{
    TODO
}
